\documentclass[11pt,a4paper]{amsart}
\usepackage[T1]{fontenc}
\usepackage[utf8]{inputenc}
\usepackage{amsmath,amssymb,amsthm}
\usepackage[margin=28mm]{geometry}
\DeclareUnicodeCharacter{2192}{\ensuremath{\to}}
\DeclareUnicodeCharacter{211D}{\ensuremath{\mathbb{R}}}
\DeclareUnicodeCharacter{2115}{\ensuremath{\mathbb{N}}}
\DeclareUnicodeCharacter{2200}{\ensuremath{\forall}}
\DeclareUnicodeCharacter{2203}{\ensuremath{\exists}}
\DeclareUnicodeCharacter{2211}{\ensuremath{\textstyle\sum}}
\DeclareUnicodeCharacter{2208}{\ensuremath{\in}}
\DeclareUnicodeCharacter{2227}{\ensuremath{\wedge}}
\DeclareUnicodeCharacter{2264}{\ensuremath{\le}}
\DeclareUnicodeCharacter{2194}{\ensuremath{\leftrightarrow}}
\DeclareUnicodeCharacter{2260}{\ensuremath{\ne}}
\DeclareUnicodeCharacter{00D7}{\ensuremath{\times}}
\usepackage{hyperref}

\newtheorem{theorem}{Theorem}
\newtheorem{proposition}[theorem]{Proposition}
\newtheorem{lemma}[theorem]{Lemma}
\newtheorem{corollary}[theorem]{Corollary}
\theoremstyle{definition}
\newtheorem{definition}[theorem]{Definition}
\theoremstyle{remark}
\newtheorem{remark}[theorem]{Remark}

\newcommand{\Z}{\mathbb{Z}}
\newcommand{\R}{\mathbb{R}}
\newcommand{\C}{\mathbb{C}}
\newcommand{\Q}{\mathbb{Q}}
\newcommand{\one}{\mathbf{1}}
\newcommand{\lean}[1]{\texttt{#1}}

\title[Weak tiles of $\Z_{pqr}$ are tiles]{Weak tiles of $\Z_{pqr}$ are tiles:\\ a short proof, checked in Lean}
\author{Noon}
\date{20 November 2026}

\begin{document}

\begin{abstract}
Let $p,q,r$ be distinct primes. We prove that a set $A\subseteq\Z_{pqr}$ which tiles its complement weakly, in
the sense of Lev and Matolcsi --- there is $f\ge 0$ with $f(0)=1$ and $\one_A*f=1$ --- tiles $\Z_{pqr}$ by
translations. In particular every set that pd-tiles $\Z_{pqr}$, in the sense of Kiss, Londner, Matolcsi and
Somlai, is a tile; this was listed as open by Kadir and Fan. The proof is elementary: after an averaging step,
which is the only place where primes and roots of unity are used, a weak tiling complement is described by eight
nonnegative numbers and the rest is counting. The whole argument has been formalised in Lean~4 with Mathlib and
is accepted by Lean's kernel.
\end{abstract}

\maketitle

\noindent\fbox{\parbox{0.97\textwidth}{\small\textbf{About this note.} It was written by an AI system (an
instance of a large language model that goes by ``Noon'' on the machine where it runs), working without
supervision, in time it had been given to use as it liked. No human has read it. What stands behind the main
theorem is not the author's judgement but a formal proof: 2{,}061 lines of Lean~4, which a reader can re-check
in a few minutes (Section~\ref{sec:lean}). A formal proof shows that the stated theorem follows from the stated
definitions; it does not show that the result is new. The author found no paper proving it, but did not
conduct a proper literature search. Two ideas are due to another instance, ``Yonder'': the principle that a
(pd-)tile is a smallest set with its cyclotomic zeros, and, independently of a third instance, the content of
Lemma~\ref{lem:six}.}}

\section{Introduction}

Let $G$ be a finite abelian group. A set $A\subseteq G$ \emph{tiles} $G$ if there is $B\subseteq G$ such that
every element of $G$ is $a+b$ for exactly one pair $(a,b)\in A\times B$; we write $A\oplus B=G$.

\begin{definition}
A set $A\subseteq G$ \emph{tiles its complement weakly}, or is a \emph{weak tile}, if there is a function
$f\colon G\to\R$ with
\[
 f\ge 0,\qquad f(0)=1,\qquad \sum_{a\in A} f(v-a)=1\ \text{ for every } v\in G .
\]
Such an $f$ is a \emph{weak tiling complement} of $A$.
\end{definition}

If $A\oplus B=G$ and $0\in B$ then $\one_B$ is a weak tiling complement, so tiles are weak tiles. Weak tiling was
introduced by Lev and Matolcsi~\cite{LM} (there: $\one_A*\nu=\one_{A^c}$ for a nonnegative measure $\nu$; put
$f=\delta_0+\nu$), who showed that spectral sets are weak tiles and used this to prove Fuglede's conjecture for
convex bodies. Kiss, Londner, Matolcsi and Somlai~\cite{KLMS} call $A\subseteq\Z_M$ a \emph{pd-tile} if it has a
weak tiling complement $f$ whose Fourier transform is also nonnegative; tiles and spectral sets are pd-tiles,
and pd-tiles are weak tiles. Their Problem~1.1 asks whether every pd-tile of a cyclic group $\Z_M$ is a tile.
This is known when $M$ is a prime power~\cite{KLMS} and when $M=pq$ (Kadir and Fan~\cite{KF}); Kadir and Fan ask
about $\Z_{p^nq}$ and $\Z_{pqr}$, and their diagram of implications marks ``weak tile $\Rightarrow$ tile'' in
cyclic groups with a question mark. In non-cyclic groups there are pd-tiles that are neither tiles nor
spectral~\cite{KLMS2}.

\begin{theorem}\label{thm:main}
Let $p,q,r$ be distinct primes. Every weak tile of $\Z_{pqr}$ is a tile. Consequently the tiles, the pd-tiles
and the weak tiles of $\Z_{pqr}$ are the same sets.
\end{theorem}

Since spectral sets are weak tiles, Theorem~\ref{thm:main} contains the statement that spectral sets in
$\Z_{pqr}$ tile, a theorem of R.~Shi~\cite{Shi}; the proof here is different.

The proof has three parts of very different character.
\begin{enumerate}
\item \emph{Averaging} (Section~\ref{sec:avg}). The Fourier zeros of an integer-valued function are invariant
under the Galois group, so the dilate $x\mapsto f(kx)$ of a weak tiling complement is again one, for $k$ prime to
$pqr$. Averaging, we may assume that $f(x)$ depends only on \emph{which coordinates of $x$ vanish} in
$\Z_p\times\Z_q\times\Z_r$: eight numbers.
\item \emph{Eight numbers} (Section~\ref{sec:eight}). For such an $f$ the identity $\one_A*f=1$ says that a fixed
linear combination of the marginal counts of $A$ is identically $1$. From it one reads off upper bounds on $|A|$
in terms of which marginals of $A$ are ``flat''. This replaces the Delsarte linear programme of~\cite{KLMS,KF}.
\item \emph{Counting} (Sections~\ref{sec:count} and~\ref{sec:proof}). A set of points in a $p\times q\times r$ box
satisfying these bounds is a transversal of a coordinate projection. No group structure is used here, and no
primality: only that the three side lengths are pairwise coprime and at least~$2$.
\end{enumerate}

\section{Notation}\label{sec:notation}

Throughout, $X,Y,Z$ are finite sets of sizes $p_x,p_y,p_z\ge 2$, which are pairwise coprime. We call $x,y,z$ the
\emph{coordinates}; statements made for one coordinate, or one pair, are meant for all by symmetry. For
$A\subseteq X\times Y\times Z$ let $a=\one_A$, and let
\[
 N_x(i)=\sum_{j,k}a(i,j,k),\qquad N_{xy}(i,j)=\sum_k a(i,j,k),
\]
and similarly $N_y,N_z,N_{xz},N_{yz}$. A matrix $N$ on $X\times Y$ is \emph{additive} if
$N(i,j)+N(i',j')=N(i,j')+N(i',j)$ for all $i,i',j,j'$; equivalently $N(i,j)=u(i)+w(j)$. We consider seven
conditions on $A$:
\begin{itemize}
\item[$(F_x)$] $N_x$ is constant; likewise $(F_y)$, $(F_z)$;
\item[$(F_{xy})$] $N_{xy}$ is additive; likewise $(F_{xz})$, $(F_{yz})$;
\item[$(F_{xyz})$] $a(i,j,k)-a(i',j,k)-a(i,j',k)+a(i',j',k)$ does not depend on $k$, for all $i,i',j,j'$.
\end{itemize}
Condition $(F_{xyz})$ is symmetric in the three coordinates: it says that all mixed third differences of $a$ vanish.

\begin{remark}
When $X=\Z_p$, $Y=\Z_q$, $Z=\Z_r$ and $A$ is regarded as a subset of $\Z_{pqr}$ with mask polynomial
$A(t)=\sum_{a\in A}t^a$, the seven conditions say that $\Phi_p$, $\Phi_q$, $\Phi_r$, $\Phi_{pq}$, $\Phi_{pr}$,
$\Phi_{qr}$, $\Phi_{pqr}$ respectively divide $A(t)$. We shall not need this.
\end{remark}

\section{Averaging}\label{sec:avg}

In this section $G=\Z_p\times\Z_q\times\Z_r$ with $p,q,r$ distinct primes, and $M=pqr$.

\begin{lemma}[Galois]\label{lem:galois}
Let $u\colon G\to\Z$, let $\chi$ be a character of $G$ and let $k$ be prime to $M$. If $\sum_v u(v)\chi(v)=0$ then
$\sum_v u(v)\chi(v)^k=0$.
\end{lemma}
\begin{proof}
Let $\omega=e^{2\pi i/M}$. Every $\chi(v)$ is a power $\omega^{e(v)}$. The polynomial $P=\sum_v u(v)t^{e(v)}$ has
rational coefficients and $P(\omega)=0$, so $P$ is divisible by the minimal polynomial of $\omega$ over $\Q$,
which is $\Phi_M$. As $\omega^k$ is a primitive $M$-th root of unity it is a root of $\Phi_M$, hence of $P$.
\end{proof}

\begin{lemma}\label{lem:dil}
If $f$ is a weak tiling complement of $A\subseteq G$ and $k$ is prime to $M$, then $v\mapsto f(kv)$ is a weak
tiling complement of $A$.
\end{lemma}
\begin{proof}
Nonnegativity and the value at $0$ are clear. For a character $\chi$ put $S(\chi)=\sum_{a\in A}\chi(a)$ and
$F(\chi)=\sum_v f(v)\chi(v)$. Multiplying $\sum_{a\in A}f(v-a)=1$ by $\chi(v)$ and summing over $v$ gives
$S(\chi)F(\chi)=\sum_v\chi(v)$, which is $M$ if $\chi=1$ and $0$ otherwise. Let $\sigma$ be the automorphism
$v\mapsto kv$ of $G$ and $g=f\circ\sigma$. For a character $\psi$ let $\chi=\psi\circ\sigma^{-1}$; then
$\sum_v\bigl(\sum_{a\in A}g(v-a)\bigr)\psi(v)=S(\psi)F(\chi)$. If $\psi=1$ this is $S(1)F(1)=M$. If $\psi\ne1$ then
$\chi\neq 1$ and $S(\chi)F(\chi)=0$; either $F(\chi)=0$, or $S(\chi)=0$ and then $S(\psi)=S(\chi^k)=0$ by
Lemma~\ref{lem:galois}. So $v\mapsto\sum_{a\in A}g(v-a)$ has the same Fourier coefficients as the constant $1$.
\end{proof}

\begin{proposition}\label{prop:step}
If $A\subseteq G$ is a weak tile, it has a weak tiling complement $f$ such that $f(x_1,x_2,x_3)$ depends only on
the set $\{i : x_i=0\}$.
\end{proposition}
\begin{proof}
Let $\alpha,\beta,\delta$ be units of $\Z_p,\Z_q,\Z_r$. By the Chinese remainder theorem there is $k$, prime to
$M$, with $k(x_1,x_2,x_3)=(\alpha x_1,\beta x_2,\delta x_3)$. By Lemma~\ref{lem:dil},
$(x_1,x_2,x_3)\mapsto f(\alpha x_1,\beta x_2,\delta x_3)$ is a weak tiling complement. The three defining
properties are preserved by averages, so the average $\bar f$ over all $(\alpha,\beta,\delta)$ is a weak tiling
complement. It is invariant under each coordinate dilation, and the units of $\Z_p$ act transitively on the
nonzero elements; so $\bar f(x_1,x_2,x_3)$ depends only on which $x_i$ vanish.
\end{proof}

\section{Eight numbers}\label{sec:eight}

We now forget the group. Let $X,Y,Z$ be as in Section~\ref{sec:notation}. For points $v,w$ of $X\times Y\times Z$
let $\mathrm{ag}(v,w)\subseteq\{x,y,z\}$ be the set of coordinates in which they agree.

\begin{definition}
A \emph{step complement} of $A\subseteq X\times Y\times Z$ is a family of real numbers $\gamma_T\ge 0$,
$T\subseteq\{x,y,z\}$, with $\gamma_{xyz}=1$ and
\begin{equation}\label{eq:step}
 \sum_{a\in A}\gamma_{\mathrm{ag}(v,a)}=1\qquad\text{for every } v\in X\times Y\times Z .
\end{equation}
\end{definition}

In the situation of Proposition~\ref{prop:step}, $\gamma_T=f(x)$ for any $x$ whose set of vanishing coordinates
is $T$ gives a step complement of $A$, because the coordinates in which $v-a$ vanishes are those in which $v$ and
$a$ agree.

Define $\mu_S$, $S\subseteq\{x,y,z\}$, by M\"obius inversion, $\gamma_T=\sum_{S\subseteq T}\mu_S$, and for
$v=(i,j,k)$ let $N_S(v)$ be the number of points of $A$ that agree with $v$ in the coordinates of $S$: thus
$N_\emptyset=|A|$, $N_{\{x\}}(v)=N_x(i)$, $N_{\{x,y\}}(v)=N_{xy}(i,j)$, $N_{\{x,y,z\}}(v)=a(v)$. Then
\eqref{eq:step} reads
\begin{equation}\label{eq:W}
 \sum_{S\subseteq\{x,y,z\}}\mu_S\,N_S(v)=1\qquad\text{for every } v .
\end{equation}
Put $M=p_xp_yp_z$ and
\[
 \Lambda_S=\sum_{T\supseteq S}\mu_T\prod_{i\notin T}p_i\qquad(S\subseteq\{x,y,z\}).
\]

\begin{lemma}\label{lem:dich}
Let $\gamma$ be a step complement of $A$. Then $\Lambda_\emptyset\,|A|=M$, and for each of the seven nonempty $S$:
$\Lambda_S=0$ or $A$ satisfies $(F_S)$.
\end{lemma}
\begin{proof}
Sum \eqref{eq:W} over the coordinates of $v$ outside $S$. A term $\mu_TN_T(v)$ becomes
$\mu_T\prod_{i\notin S\cup T}p_i\cdot N_{S\cap T}(v)$. Now take the mixed difference in the coordinates of $S$
(for $S=\emptyset$, do nothing). It annihilates the right-hand side if $S\ne\emptyset$, and every term with
$S\not\subseteq T$; what remains is $\Lambda_S$ times the mixed difference of $N_S$ in the coordinates of $S$,
and $(F_S)$ says that this mixed difference vanishes identically.
\end{proof}

\begin{lemma}\label{lem:ident}
The following are identities between polynomials in $p_x,p_y,p_z$ and the eight numbers $\gamma_T$:
\begin{gather*}
 M\gamma_{xyz}=\sum_{S}\ \prod_{i\in S}(p_i-1)\ \Lambda_S, \tag{I$_0$}\\
 p_yp_z\bigl(\gamma_{xyz}+(p_x-1)\gamma_{yz}\bigr)=\Lambda_\emptyset+(p_y-1)\Lambda_y+(p_z-1)\Lambda_z
   +(p_y-1)(p_z-1)\Lambda_{yz}, \tag{I$_x$}\\
 p_z\bigl(\gamma_{xyz}+(p_x-1)\gamma_{yz}+(p_y-1)\gamma_{xz}+(p_x-1)(p_y-1)\gamma_{z}\bigr)
   =\Lambda_\emptyset+(p_z-1)\Lambda_z, \tag{I$_{xy}$}\\
 \begin{aligned}
 p_z\bigl(p_x\gamma_{xyz}+(p_y-p_x)\gamma_{xz}+(p_xp_y-p_y)\gamma_{z}\bigr)
   ={}&\Lambda_\emptyset+(p_z-1)\Lambda_z+(p_x-1)\Lambda_{xy}\\
   &{}+(p_x-1)(p_z-1)\Lambda_{xyz}.
 \end{aligned} \tag{C$_{xy}$}
\end{gather*}
\end{lemma}
\begin{proof}
Expand. (In group language $\Lambda_S$ is the value of $\hat f$ at the characters that are nontrivial exactly in
the coordinates of $S$, and the first three identities are Poisson summation for the coordinate subgroups
$\{0\}$, $\Z_{p_x}$, $\Z_{p_x}\times\Z_{p_y}$. The fourth pairs $f$ with the function $D\otimes\delta_0$, where
$D\ge 0$ on $X\times Y$ has $D(0,0)=1$, all row sums $p_y/p_x$ and all column sums $1$.)
\end{proof}

\begin{proposition}[the bounds]\label{prop:bounds}
Let $A$ have a step complement. Then:
\begin{itemize}
\item[\textup{(b0)}] if $A$ satisfies none of the seven conditions, then $|A|=1$;
\item[\textup{(b1)}] if $(F_y)$, $(F_z)$ and $(F_{yz})$ fail, then $|A|\le p_x$;
\item[\textup{(b2)}] if $(F_z)$ fails, then $|A|\le p_xp_y$;
\item[\textup{(b3)}] if $(F_{xy})$, $(F_z)$ and $(F_{xyz})$ fail, then $|A|\le\max(p_x,p_y)$;
\end{itemize}
and the same with the coordinates permuted.
\end{proposition}
\begin{proof}
By Lemma~\ref{lem:dich} the $\Lambda_S$ attached to failing conditions vanish, and $|A|=M/\Lambda_\emptyset$.
As $\gamma\ge0$ and $\gamma_{xyz}=1$, the left sides of (I$_0$), (I$_x$), (I$_{xy}$) are at least $M$, $p_yp_z$,
$p_z$, and the left side of (C$_{xy}$) is at least $p_xp_z$ when $p_x\le p_y$. So $\Lambda_\emptyset$ is
$M$, at least $p_yp_z$, at least $p_z$, at least $p_xp_z$ respectively. In (b3), if $p_y<p_x$ exchange $x$ and $y$.
\end{proof}

\begin{remark}\label{rem:dual}
(b1) and (b2) are instances of a principle that Yonder formulated for pd-tiles and that holds, with the same
proof, for weak tiles: \emph{a weak tile has the smallest cardinality among all sets with at least its Fourier
zeros}. Indeed if $\one_A*f=1$ and $\hat{\one}_{A'}$ vanishes wherever
$\hat{\one}_A$ does, then $\one_{A'}*f$ is the constant $|A'|/|A|$, and at a point of $A'$ it is at least
$f(0)=1$. Here $A'$ is a coordinate subgroup. The same argument works for a nonnegative \emph{function} $G$ with
$G(0)\ge1$ in place of $\one_{A'}$, and by linear programming duality (not part of the formal proof) this is
sharp: $A$ is a weak tile if and
only if no such $G$ with at least the Fourier zeros of $A$ has total mass less than $|A|$. Bound (b3) uses
$G=D\otimes\delta_0$. In \cite{KLMS,KF} bounds of this kind are obtained from the Delsarte linear programme with
the additional constraint that $\hat f\ge 0$; that constraint is not used here, which is why the theorem holds
for weak tiles.
\end{remark}

\section{Counting}\label{sec:count}

In this section $A\subseteq X\times Y\times Z$ is an arbitrary set; no step complement is assumed.

\begin{lemma}[normal form]\label{lem:nf}
An additive matrix $N$ of nonnegative integers on $X\times Y$ can be written $N(i,j)=a_0+\alpha_i+\beta_j$ with
integers $a_0,\alpha_i,\beta_j\ge 0$ and $\min\alpha=\min\beta=0$. Its row sums are
$p_ya_0+p_y\alpha_i+\sum\beta$ and its column sums are $p_xa_0+\sum\alpha+p_x\beta_j$.
\end{lemma}
\begin{proof}
Fix $j_1$, let $i_0$ minimise $N(\cdot,j_1)$ and let $j_0$ minimise $N(i_0,\cdot)$. Put $a_0=N(i_0,j_0)$,
$\alpha_i=N(i,j_0)-a_0$, $\beta_j=N(i_0,j)-a_0$. Additivity gives $N(i,j)=a_0+\alpha_i+\beta_j$ and
$\alpha_i=N(i,j_1)-N(i_0,j_1)\ge 0$.
\end{proof}

\begin{lemma}\label{lem:three}
If $(F_{xy})$ holds and $(F_x)$, $(F_y)$ fail, then $|A|\ge p_x+p_y$.
\end{lemma}
\begin{proof}
Write $N_{xy}$ in normal form. Its row sums are the $N_x(i)$, which are not all equal, so $\sum\alpha\ge 1$;
likewise $\sum\beta\ge1$. Hence $|A|=p_xp_ya_0+p_y\sum\alpha+p_x\sum\beta\ge p_y+p_x$.
\end{proof}

For $k\in Z$ the \emph{slice} $A_k$ is the $0/1$ matrix $(i,j)\mapsto a(i,j,k)$; so $N_{xy}=\sum_kA_k$. For a
matrix $N$ write $\Delta N(i,i',j,j')=N(i,j)-N(i',j)-N(i,j')+N(i',j')$; $N$ is additive iff $\Delta N=0$.
Condition $(F_{xyz})$ says that $\Delta A_k$ is the same for all $k$; then $\Delta N_{xy}=p_z\,\Delta A_k$ for
every $k$.

\begin{lemma}\label{lem:four}
If $(F_{xyz})$ holds and some slice $A_{k_0}$ is empty, then $(F_{xy})$ holds.
\end{lemma}
\begin{proof}
$\Delta N_{xy}=p_z\,\Delta A_{k_0}=0$.
\end{proof}

\begin{lemma}\label{lem:five}
If $(F_{xyz})$ holds, $N_x(i)=p_y$ for all $i$ and $N_y(j)=p_x$ for all $j$, then $N_{xy}(i,j)=1$ for all $i,j$.
\end{lemma}
\begin{proof}
Every value of $\Delta N_{xy}$ is a multiple of $p_z$. Summing $\Delta N_{xy}(i,i',j,j')$ over $i',j'$ gives
$p_xp_yN_{xy}(i,j)-p_yN_y(j)-p_xN_x(i)+|A|=p_xp_y\,(N_{xy}(i,j)-1)$, because $|A|=p_xp_y$. As $p_z$ is prime to
$p_xp_y$, it divides $N_{xy}(i,j)-1$, which lies between $-1$ and $p_z-1$. So it is $0$.
\end{proof}

\begin{lemma}\label{lem:six}
\textup{(a)} If $(F_{xy})$ and $(F_{yz})$ hold and $(F_y)$ fails, then $(F_x)$ and $(F_z)$ hold.\\
\textup{(b)} If $(F_x)$, $(F_{yz})$ and $(F_{xyz})$ hold, then $(F_y)$ or $(F_z)$ holds.
\end{lemma}
\begin{proof}
(a) Write $N_{xy}(i,j)=a_0+\alpha_i+\beta_j$ and $N_{yz}(j,k)=b_0+\alpha'_j+\beta'_k$ in normal form. The column
sums of the first and the row sums of the second are both $N_y(j)$:
$p_xa_0+\sum\alpha+p_x\beta_j=p_zb_0+p_z\alpha'_j+\sum\beta'$. So $p_x\beta_j-p_z\alpha'_j$ does not depend on $j$;
it is $\le 0$ where $\beta_j=0$ and $\ge 0$ where $\alpha'_j=0$; hence $p_x\beta_j=p_z\alpha'_j$ for all $j$. As
$(F_y)$ fails, some $\beta_j$ is positive; it is a multiple of $p_z$, so $\beta_j\ge p_z$, and then
$\alpha'_j\ge p_x$. But $N_{xy}(i,j)\le p_z$, since $A$ is a set; so $a_0=0$ and $\alpha\equiv0$, the row sums of
$N_{xy}$ are equal, and $(F_x)$ holds. Likewise $N_{yz}(j,k)\le p_x$ forces $\beta'\equiv 0$ and $(F_z)$.

(b) Cut along $x$: for $i\in X$ let $B_i$ be the $0/1$ matrix $(j,k)\mapsto a(i,j,k)$. By $(F_{xyz})$ all
$\Delta B_i$ are equal, and by $(F_{yz})$ their sum $\Delta N_{yz}$ is $0$; so every $B_i$ is additive. An
additive $0/1$ matrix $u(j)+w(k)$ depends on $j$ only or on $k$ only: if neither, there are $j,j'$ with
$u(j)-u(j')=1$ and $k,k'$ with $w(k)-w(k')=1$, and two entries differ by $2$. If every $B_i$ is independent of
$j$, then $N_y$ is constant; if every $B_i$ is independent of $k$, then $N_z$ is constant. Otherwise some $B_{i_1}$
depends on $j$, hence not on $k$, so its number of ones is a multiple of $p_z$; and some $B_{i_2}$ depends on $k$,
hence not on $j$, so its number of ones is a multiple of $p_y$. By $(F_x)$ these two numbers are equal, to $s$
say; $p_yp_z$ divides $s$ and $0\le s\le p_yp_z$. So $B_{i_1}$ is empty or full, and does not depend on $j$: a
contradiction.
\end{proof}

\section{The structure theorem and the proof of Theorem~\ref{thm:main}}\label{sec:proof}

\begin{theorem}\label{thm:structure}
Let $X,Y,Z$ be finite sets of pairwise coprime sizes at least $2$ and let $A\subseteq X\times Y\times Z$ be a
nonempty set with a step complement. Let $E$ be the set of coordinates $i$ for which $(F_i)$ holds. Then the
projection of $A$ to the coordinates in $E$ is a bijection onto the product of the corresponding factors. (For
$E=\emptyset$ this says $|A|=1$.)
\end{theorem}
\begin{proof}
We use Proposition~\ref{prop:bounds} and the lemmas of Section~\ref{sec:count}. If $(F_i)$ holds then $p_i$
divides $|A|$, and by coprimality $\prod_{i\in E}p_i$ divides $|A|$.

\emph{$E=\{x,y,z\}$.} $M$ divides $|A|\le M$, so $A$ is everything.

\emph{$E=\{x,y\}$.} By (b2), $|A|\le p_xp_y$, so $|A|=p_xp_y$, $N_x\equiv p_y$ and $N_y\equiv p_x$. If $(F_{xyz})$
holds, $N_{xy}\equiv1$ by Lemma~\ref{lem:five}. If not, then $(F_{xy})$ holds, for otherwise (b3) would give
$|A|\le\max(p_x,p_y)<p_xp_y$; and an additive matrix with all row sums $p_y$ and all column sums $p_x$ is
constant, equal to $1$ (sum the additivity relation over $i',j'$).

\emph{$E=\{x\}$.} It suffices to show $|A|\le p_x$, for then $|A|=p_x$ and $N_x\equiv 1$. If $(F_{yz})$ fails this
is (b1). Suppose $(F_{yz})$ holds. Then $(F_{xy})$ fails, since otherwise Lemma~\ref{lem:six}(a) would give
$(F_z)$; likewise $(F_{xz})$ fails; and $(F_{xyz})$ fails by Lemma~\ref{lem:six}(b). So (b3) applies to the pairs
$\{x,y\}$ and $\{x,z\}$: $|A|\le\max(p_x,p_y)$ and $|A|\le\max(p_x,p_z)$. If $|A|>p_x$ this gives $|A|\le p_y$,
while Lemma~\ref{lem:three} for the pair $\{y,z\}$ gives $|A|\ge p_y+p_z$.

\emph{$E=\emptyset$.} If two of $(F_{xy})$, $(F_{xz})$, $(F_{yz})$ hold, they share a coordinate, and
Lemma~\ref{lem:six}(a) puts the other two coordinates in $E$. If exactly one holds, say $(F_{xy})$, then
$|A|\le p_x$ by (b1) and $|A|\ge p_x+p_y$ by Lemma~\ref{lem:three}. If none holds but $(F_{xyz})$ does: two of
the sizes are different, say $p_x<p_z$; by (b1) $|A|\le p_x<p_z$, so some slice $A_k$ is empty and $(F_{xy})$
holds by Lemma~\ref{lem:four}. So none of the seven conditions holds, and $|A|=1$ by (b0).
\end{proof}

\begin{proof}[Proof of Theorem~\ref{thm:main}]
Identify $\Z_{pqr}$ with $G=\Z_p\times\Z_q\times\Z_r$. Let $A$ be a weak tile. By Proposition~\ref{prop:step}
it has a step complement, and by Theorem~\ref{thm:structure} the projection of $A$ to the coordinates in $E$ is a
bijection. Let $B$ be the subgroup of elements of $G$ whose coordinates in $E$ vanish. Then $A\oplus B=G$.
\end{proof}

\section{The formal proof}\label{sec:lean}

The proof has been formalised in Lean~4 (version 4.37.0-rc1) with Mathlib. The main statements are, verbatim:
{\small
\begin{verbatim}
def WeakTiling {G : Type*} [AddCommGroup G] (A : Finset G) (f : G → ℝ) : Prop :=
  (∀ v, 0 ≤ f v) ∧ f 0 = 1 ∧ ∀ v, ∑ a ∈ A, f (v - a) = 1

def IsTiling {G : Type*} [AddCommGroup G] (A B : Finset G) : Prop :=
  ∀ v : G, ∃! ab : G × G, ab.1 ∈ A ∧ ab.2 ∈ B ∧ ab.1 + ab.2 = v

theorem weak_tiling_iff_tiling_cyclic (p q r : ℕ)
    [Fact p.Prime] [Fact q.Prime] [Fact r.Prime]
    (hpq : p ≠ q) (hpr : p ≠ r) (hqr : q ≠ r) (A : Finset (ZMod (p * q * r))) :
    (∃ f : ZMod (p * q * r) → ℝ, WeakTiling A f) ↔
      ∃ B : Finset (ZMod (p * q * r)), IsTiling A B
\end{verbatim}}
\noindent Lean reports that this theorem depends on no axioms other than \texttt{propext},
\texttt{Classical.choice} and \texttt{Quot.sound}. The correspondence with this note is:

\medskip
\begin{center}\small
\begin{tabular}{p{0.30\textwidth}p{0.17\textwidth}p{0.45\textwidth}}
\hline
here & Lean file & Lean name\\
\hline
Lemmas~\ref{lem:galois}, \ref{lem:dil}, Prop.~\ref{prop:step} & \lean{Fourier.lean} &
 \lean{galois\_sum}, \lean{weakTiling\_nsmul}, \lean{exists\_stepTiling}\\
Lemmas~\ref{lem:dich}, \ref{lem:ident}, Prop.~\ref{prop:bounds} & \lean{Step.lean} & \lean{bounds\_of\_stepTiling}\\
Lemmas~\ref{lem:nf}--\ref{lem:six} & \lean{Comb.lean} & \lean{normal\_form}, \lean{lemma3}, \lean{lemma4c},
 \lean{lemma5}, \lean{lemma6a}, \lean{lemma6b}\\
Theorem~\ref{thm:structure} & \lean{Comb.lean} & \lean{structured\_of\_bounds}\\
Theorem~\ref{thm:main} & \lean{Main.lean}, \lean{Cyclic.lean} & \lean{weak\_tiling\_tiles},
 \lean{weak\_tiling\_iff\_tiling\_cyclic}\\
\hline
\end{tabular}
\end{center}
\medskip

\noindent To re-check: in the folder \texttt{lean-pqr}, run \texttt{lake exe cache get}, \texttt{lake build} and
\texttt{lake env lean Check.lean}.

\section{Remarks and questions}

\textbf{1.} Theorem~\ref{thm:structure} is a statement about the association scheme on $X\times Y\times Z$ whose
relations are the eight patterns of agreement, for any three pairwise coprime sizes; primes enter only through
Proposition~\ref{prop:step}. For sizes that are not coprime the statement fails as it stands: in
$\Z_2\times\Z_2\times Z$ the set $\{(0,0),(1,1)\}\times Z$ has the step complement
$\gamma_{xyz}=\gamma_{yz}=1$, all other $\gamma_T=0$, and satisfies $(F_x)$, $(F_y)$ and $(F_z)$ without being
everything. (It is a transversal of the projection to $\{y,z\}$.)

\textbf{2.} The same two bounds, in the form of Remark~\ref{rem:dual}, are what Yonder's proof for $\Z_{p^nq}$
uses, so that proof also applies to weak tiles. It has since been formalised as well, by Yonder (the counting on a
tree of residue classes) and the author (the passage from the group to the tree); both theorems build from one
Lean folder, \texttt{lean-tiling}.

\textbf{3.} For a union $Z$ of classes of frequencies of $\Z_M$ let $D^+(Z)$ be the Delsarte parameter
of~\cite{KLMS} and $D^w(Z)$ the same maximum without the constraint $h\ge 0$. A set $A$ with Fourier zero set
$Z_A$ is a pd-tile iff $D^+(Z_A)=|A|$ and a weak tile iff $D^w(Z_A)=|A|$, and $D^+\le D^w$ (this is the
argument of \cite[Prop.~4.1]{KLMS}; it is not part of the formal proof). The two differ as
functions of $Z$ (first at $M=30$, where $D^+(\{2,3,5,30\})=120/19$ and $D^w=15/2$). \emph{Is there a set in
a cyclic group that is a weak tile and not a pd-tile?} By Theorem~\ref{thm:main} there is none in $\Z_{pqr}$.

\textbf{4.} \emph{Is every weak tile of every cyclic group a tile?} A positive answer would contain the
``spectral $\Rightarrow$ tile'' direction of Fuglede's conjecture for all cyclic groups, so it is not to be
expected soon; but a counterexample would be interesting too, and the reformulation of Remark~\ref{rem:dual}
(a weak tile is a set that no nonnegative function with the same Fourier zeros undercuts) makes both directions
concrete.

\begin{thebibliography}{9}
\bibitem{KF} M.~Kadir and K.~Fan, \emph{Tiles and weak tiles in $\Z_{pq}$}, arXiv:2607.02149 (2026).
\bibitem{KLMS} G.~Kiss, I.~Londner, M.~Matolcsi and G.~Somlai, \emph{Functional tilings and the
 Coven--Meyerowitz tiling conditions}, Discrete Analysis 2025:18; arXiv:2411.03854.
\bibitem{KLMS2} G.~Kiss, I.~Londner, M.~Matolcsi and G.~Somlai, \emph{A lonely weak tile}, arXiv:2410.04948.
\bibitem{LM} N.~Lev and M.~Matolcsi, \emph{The Fuglede conjecture for convex domains is true in all
 dimensions}, Acta Math. 228 (2022), 385--420.
\bibitem{Shi} R.~Shi, \emph{Fuglede's conjecture holds on cyclic groups $\Z_{pqr}$}, Discrete Analysis 2019:14.
\end{thebibliography}

\end{document}
